% Intended LaTeX compiler: pdflatex
\documentclass[a4paper,12pt]{article}
\usepackage[utf8]{inputenc}
\usepackage[T1]{fontenc}
\usepackage{graphicx}
\usepackage{longtable}
\usepackage{wrapfig}
\usepackage{rotating}
\usepackage[normalem]{ulem}
\usepackage{amsmath}
\usepackage{amssymb}
\usepackage{capt-of}
\usepackage{hyperref}
\usepackage{\string~/.emacs.d/common}
\author{mahmood}
\date{\today}
\title{closed form}
\hypersetup{
 pdfauthor={mahmood},
 pdftitle={closed form},
 pdfkeywords={},
 pdfsubject={},
 pdfcreator={Emacs 30.0.50 (Org mode 9.6)}, 
 pdflang={English}}
\begin{document}

\maketitle
to solve a , we find a \textbf{closed form} for it, a closed form for \(T(n)\) is an equation that defines \(T(n)\) using an expression that does \textbf{not} involve \(T\)\\[0pt]
there is no single method that works for all, we check for patterns and use substitution\\[0pt]
to get a better understanding, we look at examples\\[0pt]
\begin{my_example}
consider the following function that returns the factorial of a given number\\[0pt]
\begin{lstlisting}[language=Python,label= ,caption= ,captionpos=b,numbers=none]
def fac(n):
  if n == 1: return 1
  else: return n * fac(n-1)
print([fac(i) for i in range(1,10)])
\end{lstlisting}
for \texttt{fac(n)}, how many times is \texttt{fac} called?\\[0pt]
\[ T(n) = T(n-1) + C_1 \]\\[0pt]
\(C_1\) represents the total time of all the operations that take constant time that the function does on every iteration\\[0pt]
now this is obviously not what we wanna arrive at, we need to arrive at an expression that doesnt have \(T\) in it\\[0pt]
we can substitute in \(n - 1\) because \(T\) is a function that takes any positive integer\\[0pt]
\begin{align*}
  T(n) &= \underbrace{T(n-2) + C_1}_{T(n-1)} + C_1\\
  T(n) &= T(n-3) + 3C_1
\end{align*}
now you probably already see the pattern here, on every iteration the function takes constant time \(C_1\) to run, until it reaches the last iteration \(k\) at which point it we would've \(k\) times \(C_1\), e.g.:\\[0pt]
\[ T(n) \leq T(n-k) + kC_1 \]\\[0pt]

now we know that the total number of times the function calls itself depends on the initial value of \(n\), in this case that number would be \(n - 1\), as in, this function calls itself \(n - 1\) times\\[0pt]
we substitute \(k = n - 1\) to get the total time this function takes to run\\[0pt]
\begin{align*}
 T(n) &= T(n-(n-1)) + (n-1)C_1\\
      &= T(1) + (n-1)C_1\\
      &= C_1 + (n-1)C_1 &\text{because } C_1 \geq T(1)\\
      &= nC_1
\end{align*}
so we conclude that\\[0pt]
\[
  \Theta(n) = nC_1 \implies \Theta(n) = n
\]\\[0pt]
\end{my_example}
\begin{my_example}
consider the following function that returns the maximum value in a given array\\[0pt]
\begin{lstlisting}[language=Python,label= ,caption= ,captionpos=b,numbers=none]
def mymax(arr, l, r):
  if l == r:
    return arr[l]
  return max(arr[l], mymax(arr, l+1, r))
arr = [3, 10, 81, 30, 58, 0, -10]
print(mymax(arr, 0, len(arr)-1))
\end{lstlisting}
we need to find a recurrence for this function so we can analyze its time complexity\\[0pt]
first thing we notice is on every iteration we have a constant time (excluding the recursive call), we call it \(C\)\\[0pt]
we can also notice that the variable "\(r\)" never changes, so our recurrence should only depends on \(l\)\\[0pt]
now we can guess that the recurrence equation is:\\[0pt]
\[ T(l) = T(l-1) + C \]\\[0pt]
you might be confused (hopefully not) that we wrote \(l-1\) even though in the recursive call its \(l+1\), this is because we dont care about the value of the variable itself but rather about how many times the variable causes the function to call itself\\[0pt]
now if we had written \(l+1\) this would suggest that given the value \(l\) the function runs \(l+1\) times which is not correct\\[0pt]
and so this equals to \(\Theta(n)\) (using the same proof as the previous problem)\\[0pt]
\end{my_example}
\begin{my_example}
consider the following function that returns the fibonacci number at the given index\\[0pt]
\begin{lstlisting}[language=Python,label= ,caption= ,captionpos=b,numbers=none]
def fib(n):
  if n in [1, 2]: return 1
  else: return fib(n-1) + fib(n-2)
print([fib(i) for i in range(1,10)])
\end{lstlisting}
for \texttt{fib(n)}, how many times is \texttt{fib} called?\\[0pt]
\begin{align*}
  T(1) &= 1\\
  T(2) &= 1\\
  T(3) &= T(1) + T(2) + 1 = 1 + 1 + 1 = 3\\
  T(4) &= T(3) + T(2) + 1 = 3 + 1 + 1 = 5\\
  T(5) &= T(4) + T(3) + 1 = 5 + 3 + 1 = 9\\
  T(n) &= T(n-1) + T(n-2) + 1 = 2^n + 1
\end{align*}
\end{my_example}
\end{document}